% =========================================================
\section{Modos de falla y niveles de detección}
\label{sec:fallas}
% =========================================================

Un sistema de automatización no se evalúa únicamente por su
comportamiento en condiciones nominales. Durante el desarrollo se
presentaron fallas de distinta naturaleza, y el criterio que se adoptó
para juzgarlas no fue su ausencia --- inevitable en un sistema que
integra tres instrumentos independientes --- sino la capacidad del
sistema para dejar constancia de ellas en lugar de aparentar
normalidad.

Esta sección documenta las anomalías observadas en sesiones reales,
ordenadas según el nivel al que fue posible detectarlas: durante la
ejecución, mediante la verificación posterior de los datos, mediante la
auditoría del código y mediante la provocación deliberada de la
condición frente al instrumento. Se cierra examinando con el mismo criterio a
la propia herramienta de verificación. Todos los casos provienen de
registros conservados y pueden reconstruirse a partir de ellos.

\subsection{Falla detectada durante la ejecución: pérdida de la lectura
de temperatura}
\label{sec:falla_temperatura}

La sesión \texttt{20260825\_192627}, descrita en la
\cref{sec:sesion_referencia}, quedó partida en dos series por una
reprogramación del microcontrolador ESP32 a mitad de sesión. Ese corte,
que allí sirvió para contrastar dos configuraciones del láser, tuvo
además una consecuencia sobre el registro de temperatura.

En la primera serie (\texttt{m0002}--\texttt{m0006}) la lectura de
temperatura descendió de forma monótona de \SI{37.27}{\celsius} a
\SI{36.97}{\celsius}, consistente con el enfriamiento de la muestra, y
las cinco capturas se registraron con \texttt{error\_flag = 0}.

Entre esa serie y la siguiente transcurrieron doce minutos, durante los
cuales se reprogramó el firmware del ESP32 y se modificó el \textit{EO
delay} del láser. Al reprogramarse, el microcontrolador se reinició y el
hilo de lectura de temperatura terminó sin que la aplicación pudiera
reactivarlo: el método de reconexión disponible verifica la frescura de
la última lectura, pero no reinicia el hilo, operación que en esa
versión requería intervención desde la interfaz. A partir de
\texttt{m0007} el campo \texttt{temperatura} conserva el valor
\SI{36.06}{\celsius} en las once capturas restantes, y el campo
\texttt{error\_flag} se establece en 1.

El registro de la condición fue correcto: las once capturas quedaron
marcadas con \texttt{error\_flag = 1} y la interfaz mostró el indicador
de advertencia mientras la secuencia continuaba. La secuencia no se
interrumpió porque la señal fotoacústica adquirida por el osciloscopio
no depende del módulo de temperatura; lo que se perdió fue la variable
de contexto, no la medición. La decisión de no abortar fue del
operador.

El valor escrito en el campo de temperatura, en cambio, no fue
correcto, y durante varios meses se documentó como si lo fuera. Las
once filas conservan \SI{36.06}{\celsius}, que es la última lectura
válida anterior al corte: un número plausible, dentro del rango
esperado, sin nada que lo distinga de una medición efectiva salvo la
bandera de su propia fila. La \cref{sec:falla_fantasma} describe qué lo
producía y por qué la explicación que se le había atribuido era falsa.

La limitación que este episodio expone no es del registro sino de la
recuperación: el sistema supo que había perdido la lectura y lo dijo,
pero no supo recobrarla por sí mismo. Esa carencia se corrigió el 8 de
septiembre de 2026 y su verificación contra hardware se reporta en la
\cref{sec:val_instrumento}.

\begin{table}[H]
  \centering
  \caption{Registro de la sesión \texttt{20260825\_192627}. Se muestran
           las capturas representativas de cada serie y la transición
           entre ellas.}
  \label{tab:sesion_25ago}
  \begin{tabular}{llrcl}
    \toprule
    \textbf{Captura} & \textbf{Modo} & \textbf{$T$ [\si{\celsius}]} &
    \textbf{\texttt{error\_flag}} & \textbf{\texttt{error\_desc}} \\
    \midrule
    \texttt{m0001} & manual & 37.39 & 1 & --- \\
    \texttt{m0002} & tiempo & 37.27 & 0 & --- \\
    \texttt{m0006} & tiempo & 36.97 & 0 & --- \\
    \midrule
    \texttt{m0007} & manual & 36.06 & 1 & --- \\
    \texttt{m0008} & tiempo & 36.06 & 1 & temperatura no detectada \\
    \texttt{m0017} & tiempo & 36.06 & 1 & temperatura no detectada \\
    \bottomrule
  \end{tabular}
\end{table}

\subsubsection*{Un segundo hallazgo en el mismo registro}

La revisión posterior del archivo reveló una anomalía que no era
visible durante la sesión. De las doce capturas con
\texttt{error\_flag = 1}, solo diez llevan descripción del error.
Las dos restantes, \texttt{m0001} y \texttt{m0007}, tienen la bandera
levantada y el campo \texttt{error\_desc} vacío.

La coincidencia no es casual: son las dos únicas capturas de la sesión
guardadas en modo manual. El manejador correspondiente de la interfaz
recibía de la consulta al módulo de temperatura tanto el estado de
error como su descripción, pero solo propagaba el primero al registro.
El resultado es una fila que declara una anomalía sin decir cuál, lo
que obliga a inferirla del contexto.

El defecto no era observable durante la ejecución --- la interfaz
mostró correctamente la advertencia en ambos casos --- ni afectaba a
las capturas automáticas, que constituyen el modo de operación
habitual. Se detectó al auditar el registro completo, y se corrigió
haciendo que el manejador propague ambos campos. La corrección no
puede aplicarse retroactivamente a los datos del 25 de agosto: las dos
filas permanecen como testimonio del defecto.

Este hallazgo ilustra el argumento central de la sección desde un
ángulo inesperado. El sistema registró suficiente información para que
su propio defecto de registro resultara detectable. Una fila con
bandera levantada y sin explicación es anómala a simple vista; una fila
sin bandera habría pasado inadvertida.

\subsection{Falla detectada al provocarla con instrumento: el valor de
temperatura conservado}
\label{sec:falla_fantasma}

El caso anterior se dio por entendido durante meses. La conservación
del último valor leído se describió como una decisión de diseño,
preferible a escribir un cero que después resultaría indistinguible de
una medición válida. La verificación con el instrumento, realizada en
la sesión del 8 de septiembre de 2026, mostró que esa descripción era
falsa en su parte esencial.

La prueba consistió en desconectar físicamente el módulo de temperatura
y guardar capturas en modo manual, condición que no puede fabricarse
por simulación: \texttt{generar\_sesion\_prueba.py} produce el archivo
resultante, no el evento que lo origina. Las tres filas obtenidas en la
sesión \texttt{20260908\_184623} (\texttt{m0012}--\texttt{m0014})
llevan \texttt{error\_flag = 1}, pero también un campo de temperatura
con el valor \SI{24.25}{\celsius} y un \texttt{error\_desc} genérico
que no nombra al dispositivo ausente. No había sensor conectado: ese
número no correspondía a ninguna medición.

La inspección del código explicó el origen. La bandera de error se
derivaba de un indicador de frescura de la lectura, que es el criterio
correcto; la escritura del campo de temperatura, en cambio, estaba
condicionada a que la variable de lectura no fuera nula. Esa variable
conserva el último valor recibido y no vuelve a ser nula una vez que ha
habido una primera lectura, de modo que la rama encargada de escribir
un valor de ausencia nunca llegaba a ejecutarse. No era una decisión de
diseño: era código inalcanzable. El sistema no conservaba el último
valor por criterio, sino porque carecía de camino para dejar de
conservarlo.

La distinción importa para el argumento de esta sección. Un registro
que marca la anomalía y además escribe un número plausible en el campo
afectado no falla del todo, pero tampoco cumple: basta que un análisis
posterior filtre por el campo de temperatura sin consultar la bandera
para que el valor arrastrado entre en el cálculo como si fuera una
medición. Una bandera protege únicamente a quien la lee.

La corrección escribe \texttt{nan} cuando la lectura no es fresca, y
hace que \texttt{error\_desc} nombre el dispositivo y el puerto. Se
descartó deliberadamente el valor \num{0.0} que se había enunciado como
criterio de aceptación en la \cref{sec:val_instrumento}: cero grados
Celsius es una temperatura físicamente válida y reintroduciría la misma
ambigüedad que se pretende eliminar, mientras que \texttt{nan} no
admite lectura como dato. La verificación en la sesión
\texttt{20260908\_190614} devuelve, con el módulo desconectado,
\texttt{error\_flag = 1}, \texttt{temperatura = nan} y un
\texttt{error\_desc} que identifica al ESP32 y su puerto.

Las once filas del 25 de agosto no admiten corrección retroactiva.
Quedan, como las dos filas sin descripción de la subsección anterior,
en calidad de testimonio: su campo de temperatura no es un dato, y su
bandera es lo único que lo advierte.

\subsection{Falla detectada en la verificación posterior: duplicación
de la adquisición transferida}
\label{sec:falla_buffer}

La falla anterior fue visible mientras ocurría. La que se describe a
continuación no lo fue: transcurrió sin señal alguna en la interfaz, sin
excepción registrada y sin marca en el archivo de sesión. Se encontró al
someter los datos a una verificación independiente, semanas después.

En la sesión \texttt{20260818\_190223}, de cuatro capturas, las
mediciones \texttt{m0003} y \texttt{m0004} resultaron idénticas bit a
bit: mismo resumen MD5 sobre las mismas 25\,000 muestras. El registro,
sin embargo, les asigna parámetros temporales distintos. A
\texttt{m0003} le corresponde un \texttt{XINCR} de \SI{80}{\nano\second}
y a \texttt{m0004} uno de \SI{4}{\nano\second}: un factor de veinte en
la escala del eje temporal, declarado sobre exactamente el mismo
arreglo de datos.

La explicación es una condición de carrera entre la configuración del
instrumento y la transferencia. Se modificó la escala horizontal desde
el panel del osciloscopio, el sistema consultó y registró el preámbulo
ya actualizado, pero la consulta \texttt{CURVE?} devolvió la
adquisición anterior, tomada con la escala previa. El resultado es un
registro internamente inconsistente: metadatos nuevos describiendo
datos viejos.

\subsubsection*{Consecuencia sobre la reconstrucción}

La inconsistencia no es cosmética. La reconstrucción del eje temporal a
partir de \texttt{XINCR}, \texttt{XZERO} y \texttt{PT\_OFF} sitúa el
máximo de \texttt{m0003} en \SI{128.64}{\micro\second} y el de
\texttt{m0004} en \SI{6.43}{\micro\second}. Se trata del mismo arreglo
de muestras interpretado sobre dos ejes distintos, de modo que al menos
una de las dos reconstrucciones es falsa, y ninguna de las dos anuncia
serlo.

Ninguna bandera lo señaló. Las cuatro capturas de esa sesión llevan
\texttt{error\_flag = 1}, pero por una razón ajena: la sesión corrió sin
el módulo de temperatura conectado y el monitor mantuvo la bandera
levantada de principio a fin. La bandera estaba encendida por el motivo
equivocado, y por lo tanto no distinguía las capturas corruptas de las
sanas. Una bandera que se enciende siempre no informa nada.

\subsubsection*{El criterio que la detecta}

La herramienta de verificación compara las capturas crudas entre sí,
sin recurrir a los metadatos. El criterio se apoya en una propiedad del
fenómeno: dos adquisiciones consecutivas de una señal real nunca
coinciden muestra por muestra, porque el ruido de fondo difiere en cada
una. En esa misma sesión, los demás pares de capturas difieren en
alrededor de 24\,000 de sus 25\,000 muestras. Una coincidencia exacta no
es un resultado improbable: es imposible, y por lo tanto delata una
transferencia que no ocurrió.

Conviene subrayar que este criterio es independiente de los metadatos
que se pretende validar. Verificar la consistencia del preámbulo consigo
mismo habría resultado satisfactorio en ambas capturas, porque cada una
es internamente coherente. Lo que revela el defecto es la comparación
entre capturas, no la inspección de cada una por separado.

\subsubsection*{Un antecedente en las sesiones preliminares}

El mismo criterio, aplicado a los registros de junio de 2026 ---
anteriores a la integración del módulo de temperatura y conservados
como material de desarrollo ---, revela que el defecto no fue un
incidente aislado.

En la sesión \texttt{20260623\_171412} las cuatro primeras capturas
registran una señal de aproximadamente \SI{2.9}{\milli\volt} de
amplitud pico a pico. Las cuatro siguientes son idénticas entre sí bit
a bit, con una amplitud de \SI{0.078}{\milli\volt}. En la sesión
\texttt{20260623\_173432}, tomada dieciocho minutos después, la
repetición abarca diez capturas consecutivas, desde \texttt{m0004}
hasta \texttt{m0013}, separadas por intervalos de noventa segundos a lo
largo de doce minutos.

Podría objetarse que un registro de amplitud tan baja corresponde
simplemente a la ausencia de señal en la entrada del osciloscopio, lo
que en una sesión de desarrollo sería esperable. La objeción no explica
lo observado. Esos registros no son constantes: presentan una
desviación estándar de \SI{0.0095}{\milli\volt}, es decir, ruido de
fondo. Dos adquisiciones sucesivas de ruido separadas por noventa
segundos no pueden coincidir muestra por muestra. La coincidencia
exacta indica que no hubo adquisición, con independencia de lo que
estuviera conectado a la entrada.

La coincidencia admite además una confirmación que no depende de la
herramienta de verificación. Al incorporar los archivos de esa sesión
al sistema de control de versiones, este calculó para cada uno un
identificador derivado de su contenido: las diez capturas de
\texttt{m0004} a \texttt{m0013} recibieron el mismo, mientras que las
tres anteriores recibieron identificadores distintos entre sí. El
resultado importa por su procedencia. El verificador compara las
capturas porque fue escrito para buscar precisamente esta anomalía; el
control de versiones no sabe qué es una forma de onda ni qué se espera
de ella, y llegó al mismo agrupamiento por el solo hecho de indexar
archivos. La duplicación queda así sostenida por dos vías
independientes.

Estas sesiones no forman parte de los resultados del sistema: se
realizaron con una versión anterior del programa, sin módulo de
temperatura y con un formato de registro que aún no incluía la
descripción de errores. Se citan aquí por lo que aportan al alcance de
la falla: la duplicación de la adquisición ocurrió al menos en tres
sesiones distintas a lo largo de dos meses, y en ninguna de ellas
levantó bandera alguna. El caso del 18 de agosto no es una anomalía
puntual, sino la manifestación más breve de una condición recurrente.

\subsubsection*{Contraste: la reconstrucción cuando los metadatos son
correctos}

El caso anterior podría sugerir que la fórmula de reconstrucción es
poco confiable. La \cref{sec:reconstruccion_escalas} demuestra lo
contrario: en la sesión del 25 de agosto, dos muestreos con
resoluciones distintas convergen tras la reconstrucción al mismo
instante físico.

Los dos casos delimitan con precisión el alcance de la reconstrucción.
La fórmula es correcta y se sostiene entre escalas; lo que puede fallar
es la correspondencia entre el arreglo de muestras y los parámetros que
se le adjuntan. Ese emparejamiento no se valida solo, y es el objeto de
la comparación descrita arriba.

\subsection{Falla detectada en la auditoría del código: regresión
silenciosa en el registro de escalas}
\label{sec:falla_regresion}

Los tres casos anteriores dejaron rastro en los datos: una advertencia
en la interfaz, un valor de temperatura escrito sin sensor que lo
produjera, un par de arreglos idénticos. El cuarto no dejó ninguno.
No produjo error, no levantó bandera y no alteró ninguna medición. Su
única manifestación fue la ausencia de dos valores que el propio archivo
declaraba contener.

Las columnas \texttt{CH\_SCALE} y \texttt{HOR\_SCALE} del registro de
sesión almacenan las escalas vertical y horizontal vigentes en el
osciloscopio al momento de cada captura. Se incorporaron el 15 de junio
de 2026, junto con las consultas \texttt{\{canal\}:SCALE?} y
\texttt{HORizontal:SCAle?} que las alimentan, situadas dentro de la
rutina de lectura del preámbulo.

Ocho días después, un cambio dedicado a incorporar la selección de
longitud de registro revirtió, sobre ese mismo archivo, la totalidad de
lo que el cambio anterior había introducido: las dos consultas de
escala, la actualización de la longitud de registro en memoria, y un
método auxiliar de estimación de tiempos. El mensaje del cambio no
menciona ninguna de las tres supresiones.

\subsubsection*{Por qué una de las tres pérdidas fue visible y dos no}

La supresión del método auxiliar se manifestó de inmediato: la primera
llamada produjo una excepción por atributo inexistente, quedó registrada
como defecto y se corrigió. Las dos consultas de escala no produjeron
nada. La rutina de lectura siguió devolviendo su diccionario de
parámetros, ahora sin esas dos claves; el módulo de almacenamiento las
solicitó con un valor por omisión y escribió la cadena vacía. Ninguna
excepción, ningún registro, ninguna diferencia perceptible durante la
operación.

La distinción es instructiva. Un valor que alguien consume falla
ruidosamente cuando desaparece, porque su consumidor protesta. Un valor
que solo se escribe y nunca se lee se degrada en silencio: el encabezado
del archivo lo sigue anunciando, la columna se sigue escribiendo, y lo
único que cambia es su contenido. Entre el 23 de junio y el 31 de agosto
las diecisiete capturas del 25 de agosto y las cuatro del 18 se
registraron con ambas columnas vacías.

Vale la pena señalar un detalle del código original que descarta la
explicación más cómoda. Las dos consultas estaban protegidas por
manejadores de excepción que, ante un fallo de comunicación, asignaban
el valor \num{0.0}. Un osciloscopio que no respondiera habría producido
ceros, no celdas vacías. El vacío observado no indica un instrumento que
no contestó, sino un código que no preguntó.

\subsubsection*{Detección y corrección}

El defecto se encontró al contrastar el encabezado declarado del archivo
de sesión contra su contenido efectivo, columna por columna, y rastrear
después el historial del archivo fuente para ubicar el punto de
supresión. Ninguno de los dos niveles anteriores podía alcanzarlo: en
ejecución no hay nada que observar, y la verificación de los datos
compara capturas entre sí, no el archivo contra su propia promesa.

La corrección restituyó ambas consultas dentro de la rutina de lectura
del preámbulo, de modo que las escalas se obtengan en la misma operación
que el resto de los parámetros de la captura y no en una llamada
separada susceptible de quedar desincronizada.
Esa verificación se realizó el 8 de septiembre de 2026. En la sesión
\texttt{20260908\_184623} las once capturas registran
\texttt{CH\_SCALE} $=\num{0.001}$ y \texttt{HOR\_SCALE} $=\num{1e-5}$,
equivalentes a \SI{1}{\milli\volt} y \SI{10}{\micro\second} por
división, en coincidencia con el panel del instrumento. Ninguna celda
quedó vacía, y ninguna registró el valor \num{0.0} que habrían
producido los manejadores de excepción ante un fallo de comunicación.

\subsection{La herramienta de verificación como fuente de artefacto}
\label{sec:falla_umbral}

Los cuatro casos anteriores describen fallas del sistema de adquisición
detectadas por medios independientes. Corresponde examinar el último de
esos medios con el mismo criterio, porque también puede equivocarse, y
en esta sesión lo hizo.

Entre las magnitudes que reporta la herramienta de verificación se
encuentra el instante de inicio de la señal, definido como el primer
punto cuyo valor absoluto supera un umbral. Sobre la sesión del 25 de
agosto ese instante resultó ser de \SI{31.4}{\micro\second} en la
primera serie y de \SI{1.5}{\micro\second} en la segunda: una
diferencia de treinta microsegundos entre dos series tomadas sobre la
misma muestra con quince minutos de separación.

La discrepancia es un artefacto del criterio, no un desplazamiento de la
señal. El umbral empleado es relativo: ocho veces la desviación estándar
del ruido medido en la porción inicial de cada registro. El ruido de la
segunda serie resultó ser aproximadamente un tercio del de la primera
---\SI{0.0056}{\milli\volt} contra \SI{0.0171}{\milli\volt}---, de modo
que el mismo criterio produjo umbrales de \SI{0.045}{\milli\volt} y
\SI{0.137}{\milli\volt} respectivamente. El umbral más bajo cruza antes,
y lo hace sobre fluctuaciones que en la primera serie quedaban por
debajo del suyo.

Aplicando un umbral fijo de \SI{0.137}{\milli\volt} a las quince
capturas adquiridas en modo por tiempo, el inicio se sitúa en
\SI{31.4}{\micro\second} y \SI{31.6}{\micro\second} para las dos series.
Corroborando el resultado por una vía distinta, el instante del máximo
---que no depende de umbral alguno, pues se obtiene del índice de mayor
amplitud--- coincide en \SI{69.13}{\micro\second} y
\SI{69.16}{\micro\second}. La señal no se movió.

El umbral fijo no es, sin embargo, universalmente preferible, y el
mismo cálculo lo muestra. Aplicado a las diecisiete capturas de la
sesión, dos se apartan del valor común: la captura manual m0001 sitúa el
inicio en \SI{1.8}{\micro\second}, porque contiene una fluctuación
temprana que supera los \SI{0.137}{\milli\volt} antes del arribo real,
y m0006 lo sitúa en \SI{30.7}{\micro\second}. En m0001 el criterio
relativo acierta donde el fijo falla: como esa captura tiene un ruido
mayor, su umbral propio queda por encima de la fluctuación y no la
confunde con el inicio. Ninguno de los dos criterios es correcto en
todos los casos. El fijo hace comparables entre sí capturas de distinta
relación señal a ruido, que es lo que esta sección necesitaba, a costa
de perder la protección que el criterio relativo ofrece frente a
fluctuaciones tempranas de amplitud moderada. Se adopta el fijo por esa
razón y se declara su costo.

El episodio deja dos observaciones sobre la herramienta y una sobre el
conjunto. Sobre la herramienta: un criterio relativo produce magnitudes
que no son comparables entre capturas de distinta relación señal a
ruido, y el chequeo de deriva del inicio, que agrupa las capturas por
escala antes de compararlas, no puede detectar por construcción una
discrepancia que ocurre precisamente entre grupos de escala distinta.
Ambas limitaciones subsisten; como medida provisional, la herramienta
imprime hoy el umbral efectivo de cada captura, de modo que una
diferencia como la descrita resulta legible en la tabla en lugar de
tener que deducirse.

Sobre el conjunto: ninguno de los cuatro niveles de detección basta por
sí solo. La verificación posterior encontró lo que la ejecución no vio,
la auditoría encontró lo que la verificación no podía ver, la
provocación deliberada de la condición frente al instrumento refutó lo
que la auditoría había dado por entendido, y la verificación misma
requirió ser contrastada contra una magnitud independiente para
distinguir un hallazgo de un artefacto de su propio criterio.

La sesión del 8 de septiembre de 2026 permitió someter ese criterio a
una prueba que antes no existía: nueve capturas consecutivas del mismo
arreglo, tomadas sobre un descenso continuo de temperatura, con la
amplitud de la señal decreciendo de forma sostenida a lo largo de la
serie. El umbral fijo no resistió la prueba.

Aplicado a las nueve capturas, el criterio relativo al ruido desplaza
el instante de arribo estimado en \SI{0.39}{\micro\second} entre el
primer punto y el último. Sustituirlo por el umbral fijo adoptado
reduce ese desplazamiento a \SI{0.32}{\micro\second}: prácticamente
nada. Peor aún, la correlación del arribo estimado con la temperatura
se vuelve más estrecha, no más débil. El criterio que se había
adoptado para eliminar el artefacto lo conserva casi intacto.

La razón se encuentra al medir el mismo desplazamiento por vías que no
dependen de umbral alguno. Alineando cada captura contra la primera
mediante correlación cruzada del registro completo, la serie se
desplaza \SI{0.12}{\micro\second} en sentido \emph{contrario}: hacia
tiempos menores. El instante del máximo absoluto, que tampoco depende
de umbral, se desplaza \SI{0.52}{\micro\second} en ese mismo sentido
contrario. Las tres medidas no pueden describir el mismo fenómeno.

\begin{table}[H]
  \centering
  \caption{Desplazamiento aparente del arribo en las nueve capturas
           automáticas de la sesión \texttt{20260908\_203856}, según el
           criterio empleado para estimarlo.}
  \label{tab:umbral_amplitud}
  \begin{tabular}{lcc}
    \toprule
    Criterio & Sentido al enfriarse & Magnitud (\si{\micro\second}) \\
    \midrule
    Umbral relativo al ruido & Hacia tiempos mayores & 0.39 \\
    Umbral fijo & Hacia tiempos mayores & 0.32 \\
    Umbral proporcional a la amplitud & Hacia tiempos mayores & 0.16 \\
    Correlación cruzada del registro & Hacia tiempos menores & 0.12 \\
    Instante del máximo absoluto & Hacia tiempos menores & 0.52 \\
    \bottomrule
  \end{tabular}
\end{table}

La contradicción se resuelve al observar qué cambia efectivamente a lo
largo de la serie. La amplitud pico a pico desciende de
\SI{1.579}{\milli\volt} en el primer punto a \SI{1.322}{\milli\volt} en
el último, mientras la desviación estándar del segmento previo a la
señal sube de \SI{0.0148}{\milli\volt} a \SI{0.0185}{\milli\volt}. Un
flanco de subida de menor amplitud, con la misma forma, tiene menor
pendiente; y un umbral cualquiera, fijo o no, lo cruza más tarde
aunque la onda no se haya movido en absoluto. Lo que el criterio de
umbral mide en esta serie no es el instante de arribo: es la pérdida
de amplitud de la señal.

La confirmación es directa. Si el umbral se hace proporcional a la
amplitud de cada captura en lugar de fijo, el desplazamiento cae a
\SI{0.16}{\micro\second}, la mitad, sin ningún otro cambio en el
procedimiento.

El alcance de esta observación conviene enunciarlo sin atenuantes. La
subsección anterior adoptó el umbral fijo como corrección de un
artefacto que se había manifestado entre sesiones con niveles de ruido
distintos. Esa corrección es válida para ese caso y sigue siéndolo.
Pero el barrido por temperatura produce, por construcción, series en
las que la amplitud de la señal varía de una captura a la siguiente, y
en ese régimen ---que es el modo de operación principal del sistema---
el umbral fijo no es un criterio robusto. Ninguna versión de
\texttt{verificar\_sesion.py} empleada en este trabajo lo es. La
herramienta detecta correctamente que algo se desplaza; lo que no puede
hacer, con el criterio de umbral, es decir qué.

Quedan dos consecuencias. La primera es metodológica y acotada: un
estimador de arribo robusto frente a cambios de amplitud ---la
correlación cruzada contra una captura de referencia es el candidato
natural, y está disponible sin instrumentación adicional--- debería
sustituir al umbral en cualquier análisis que compare capturas de un
mismo barrido. La segunda no es de la herramienta sino de los datos, y
se recoge en el \cref{cap:conclusiones}: una vez descontado el
artefacto, persiste un desplazamiento residual cuyo sentido contradice
lo que predice la dependencia de la velocidad del sonido con la
temperatura.

\subsection{Anomalía detectada en auditoría, sin consecuencia
manifiesta: la frecuencia de repetición supuesta}
\label{sec:falla_laserhz}

El campo \texttt{pulsos\_estimados} del registro de sesión declara
cuántos disparos del láser quedaron integrados dentro de la ventana de
adquisición de cada punto. Se obtiene multiplicando la duración de esa
ventana por la frecuencia de repetición del láser. La duración se mide;
la frecuencia, no: es una constante escrita en el código con el valor
de \SI{10}{\hertz}.

El equipo expone su frecuencia de repetición como registro consultable
a través de la misma biblioteca de control que el sistema ya emplea
para gobernarlo. El programa nunca la lee. Mientras el láser opere a
\SI{10}{\hertz} ---como ocurrió en todas las sesiones de este
trabajo--- el número escrito es correcto, y esa es exactamente la razón
por la que la anomalía no se había advertido.

Su naturaleza difiere de las seis anteriores. Aquellas se
manifestaron: hubo columnas vacías, filas sin descripción, un valor de
temperatura que no correspondía a ninguna medición. Esta no ha
producido todavía ningún dato incorrecto. Lo que la hace anotable es su
forma, que es la misma del caso de la temperatura conservada: un campo
que el archivo presenta como resultado de una medición y que en
realidad transporta una suposición del programador. Si alguien
modificara la frecuencia de repetición del láser desde su control
externo, cosa que el equipo permite y que el sistema no vigila,
\texttt{pulsos\_estimados} seguiría reportando como si nada hubiera
cambiado, sin bandera de error y sin discrepancia visible en el
archivo. El dato correcto lo sería por coincidencia.

La corrección es de alcance reducido ---leer el registro en lugar de
suponerlo, y registrar el valor leído junto con los demás metadatos de
sesión--- y se recoge en el \cref{cap:conclusiones}. Se enuncia aquí
porque el criterio que organiza esta sección es el nivel al que una
anomalía resulta detectable, y esta solo lo fue mediante la lectura
deliberada del código en busca del origen de cada campo del archivo.
Ningún dato la habría delatado.

\subsection{Una comprobación del registro: el punto perturbado}
\label{sec:registro_corrige}

Durante la sesión del 8 de septiembre de 2026 hubo una interrupción
física del arreglo mientras la secuencia estaba en curso: se colocó una
protección sobre el láser con la adquisición corriendo. El hecho quedó
anotado en la bitácora de la sesión junto con la atribución que el
operador hizo en ese momento, señalando dos capturas concretas como
afectadas.

El análisis posterior de los archivos, realizado seis días después y
sin más elementos que los datos guardados, contradice esa atribución.
Las dos capturas señaladas no se distinguen de sus vecinas por ningún
criterio: ni su amplitud pico a pico, que sigue la tendencia
descendente de la serie, ni el instante de su máximo, que coincide con
el del resto dentro de medio microsegundo.

La captura que sí se aparta es \texttt{m0004}. Dos indicadores
independientes la señalan. Su amplitud pico a pico es
\SI{1.732}{\milli\volt}, el valor más alto de toda la sesión, en medio
de una serie que desciende de forma sostenida. Y el instante de su
máximo absoluto cae en \SI{78.08}{\micro\second}, cuando en las ocho
capturas restantes se ubica entre \SI{94.37}{\micro\second} y
\SI{94.89}{\micro\second}. No se trata de la misma señal medida con
ruido: es una señal distinta.

El episodio ilustra el argumento de este capítulo mejor que ninguna de
las anomalías anteriores, y conviene enunciarlo con precisión para no
reclamar de más. El sistema no detectó la perturbación: ningún campo
del archivo levantó bandera, y no debía hacerlo, porque nada en la
adquisición había fallado. Lo que el sistema hizo fue registrar, junto
a cada forma de onda, las condiciones bajo las que se obtuvo y la
trazabilidad suficiente para que la pregunta pudiera formularse
después. La perturbación se identificó comparando nueve capturas entre
sí, algo que solo es posible si las nueve están guardadas con sus
escalas, su temperatura y su orden.

Ese es el criterio con el que debe juzgarse un sistema de
automatización de medición: no que anticipe todo lo que puede salir
mal, sino que deje registro suficiente para reconstruirlo. En este
caso el registro corrigió al operador, que estuvo presente, sobre lo
que el operador creyó haber visto.

\subsection{Balance}
\label{sec:fallas_balance}

\begin{table}[H]
  \centering
  \caption{Anomalías documentadas, nivel al que fueron detectadas y
           estado actual.}
  \label{tab:resumen_fallas}
  \small
  \begin{tabular}{p{3.6cm}p{2.6cm}p{3.4cm}p{3.2cm}}
    \toprule
    \textbf{Anomalía} & \textbf{Nivel de detección} &
    \textbf{Evidencia} & \textbf{Estado} \\
    \midrule
    Pérdida de la lectura de temperatura & Ejecución &
    Doce capturas marcadas y advertencia en la interfaz &
    Registrada; reconexión automática implementada y verificada \\
    \addlinespace
    Valor de temperatura conservado tras la pérdida de lectura &
    Provocación con instrumento &
    Tres filas con bandera levantada y temperatura numérica sin sensor
    conectado &
    Corregida; se escribe \texttt{nan} \\
    \addlinespace
    Descripción de error omitida en el guardado manual &
    Auditoría del registro &
    Dos capturas con bandera y sin descripción & Corregida \\
    \addlinespace
    Buffer duplicado en la transferencia &
    Verificación posterior &
    Arreglos idénticos entre capturas sucesivas, en tres sesiones &
    Detectada; sin corrección en el instrumento \\
    \addlinespace
    Regresión en el registro de escalas & Auditoría del código &
    Veintiuna capturas con las columnas vacías &
    Corregida y verificada en hardware \\
    \addlinespace
    Artefacto del umbral relativo & Contraste entre métricas &
    \SI{30}{\micro\second} de diferencia aparente &
    Documentada; criterio vigente \\
    \addlinespace
    Frecuencia de repetición supuesta en el cálculo de
    \texttt{pulsos\_estimados} &
    Auditoría del código &
    Ninguna: el valor supuesto coincide con el real en todas las sesiones &
    Documentada; corrección propuesta \\
    \bottomrule
  \end{tabular}
\end{table}

Seis de las siete anomalías no eran observables durante la operación.
Una de ellas, además, se había documentado como decisión deliberada de
diseño: solo provocar la condición frente al instrumento mostró que no
lo era. La séptima no llegó siquiera a manifestarse: se detectó
leyendo el código en busca del origen de cada campo del archivo, y su
consecuencia sigue siendo hipotética.
Tres de ellas tampoco lo eran en los datos tomados uno a
uno: solo aparecieron al contrastar el archivo con lo que declaraba
contener, el código con su propia historia, y cada campo del archivo
con el lugar del código que lo produce. Ninguna produjo una
excepción ni interrumpió una secuencia. La más prolongada ---diez
capturas consecutivas sin adquisición--- transcurrió durante doce
minutos frente a un operador atento.

El sistema no evita fallar. Lo que hace es conservar información
suficiente para que sus fallas sean reconstruibles después, y esa
propiedad es la que permitió documentar esta sección. Las mejoras que se
desprenden de cada caso se recogen en el \cref{cap:conclusiones}.
